% Formats des imports et exports d'EP Manager (mod_ep).
% Compilation : latexmk -pdf formats-import.tex
\documentclass[11pt,a4paper]{article}
\usepackage[french]{babel}
\usepackage[T1]{fontenc}
\usepackage[utf8]{inputenc}
\usepackage{lmodern,microtype,booktabs,tabularx,longtable,xcolor,enumitem,hyperref,array}
\usepackage[margin=2.2cm]{geometry}
\hypersetup{colorlinks=true, linkcolor=blue!50!black, urlcolor=blue!50!black}
\newcommand{\code}[1]{\texttt{\detokenize{#1}}}
\newcommand{\col}[1]{\texttt{#1}}
\newcommand{\oblig}{\textbf{oui}}
\setlist[itemize]{nosep}
\setlist[enumerate]{nosep}
\newcolumntype{L}[1]{>{\raggedright\arraybackslash}p{#1}}

\title{EP Manager : formats des imports et des exports\\
  \large Catalogue, séances, crédits, paramétrage, export complet}
\author{Sébastien Lefebvre}
\date{Version 0.7.0 du module \code{mod_ep} -- novembre 2026}

\begin{document}
\maketitle
\tableofcontents

\section{Principes communs}

Tous les imports se font depuis le panneau \textbf{Administration} de l'activité, section
\textbf{Imports, exports et transfert}, sans passer par l'administration du site. Ils
suivent le même déroulement :
\begin{enumerate}
  \item dépôt du fichier ;
  \item \textbf{prévisualisation} : chaque ligne est lue et contrôlée, les erreurs et
    avertissements sont listés avec leur numéro de ligne, et le tableau montre ce qui sera
    créé ou mis à jour. \emph{Rien n'est encore écrit} ;
  \item \textbf{confirmation} : le fichier déposé est relu et appliqué en une seule
    transaction, sous le verrou de l'activité. Les lignes en erreur sont écartées, les
    autres appliquées. La prévisualisation expire au bout de 24~heures.
\end{enumerate}
Chaque import confirmé est inscrit au \textbf{journal des imports} (page \emph{Transférer le
paramétrage}) : date, nature, auteur, nombre d'éléments et compte rendu.

\paragraph{Fichiers CSV.} Une ligne d'en-tête, puis une ligne par élément.
\begin{itemize}
  \item Séparateur : point-virgule, virgule ou tabulation, détecté sur la ligne d'en-tête.
    Un fichier enregistré par Excel en «~CSV (séparateur : point-virgule)~» convient tel
    quel.
  \item Encodage : UTF-8, avec ou sans BOM, ou Windows-1252 (Excel sous Windows).
  \item Les colonnes sont reconnues \textbf{par leur intitulé}, dans n'importe quel ordre,
    sans tenir compte des accents, de la casse, des espaces ni de la ponctuation :
    \col{Intitulé}, \code{intitule} et \col{INTITULÉ} désignent la même colonne. Plusieurs
    intitulés sont acceptés pour chaque colonne, en français et en anglais (tableaux
    ci-dessous). Une colonne inconnue est signalée et ignorée.
  \item Un champ entre guillemets peut contenir des retours à la ligne (une rubrique de
    syllabus sur plusieurs paragraphes, par exemple).
  \item Nombres : virgule ou point décimal (\code{2,5} ou \code{2.5}) ; les espaces de
    milliers sont tolérés.
  \item Dates : \code{jj/mm/aaaa} ou \code{aaaa-mm-jj}, éventuellement suivies d'une heure
    \code{hh:mm} (ou \code{hhhmm}). Elles sont lues dans le fuseau horaire de la personne
    qui importe.
  \item Oui/non : \code{oui}, \code{o}, \code{1}, \code{x}, \code{vrai}, \code{yes}, ou
    \code{non}, \code{n}, \code{0}, \code{faux}, \code{no}.
  \item Personnes (responsables, intervenants, étudiants, validateurs) : adresse
    électronique ou numéro d'identification du compte Moodle. Une adresse portée par deux
    comptes est ambiguë et n'est rapprochée d'aucun.
\end{itemize}
Chaque page d'import affiche le tableau des colonnes reconnues et propose un
\textbf{modèle} à télécharger (la ligne d'en-tête seule). Des fichiers d'exemple complets se
trouvent dans le dossier \code{docs/exemples} du dépôt.

\section{Catalogue des EP}\label{sec:catalogue}

Page \code{import_catalog.php}, réservée à la gestion de l'activité
(\code{mod/ep:manage}). Une ligne par EP du catalogue.

\paragraph{Rapprochement.} Un EP du fichier est rapproché d'un EP existant par son
\textbf{code}, à défaut par son \textbf{intitulé} (sans tenir compte de la casse). S'il
existe, il est mis à jour : seules les colonnes présentes dans le fichier modifient ses
caractéristiques, les autres gardent leur valeur ; sans colonne \code{Code}, un EP rapproché
par son intitulé garde son code. Sinon, il est créé. Deux lignes du fichier ne peuvent pas
porter le même code. Rien n'est jamais supprimé.

\begin{longtable}{@{}L{3.2cm}L{4.6cm}L{1.4cm}L{6cm}@{}}
\toprule
Colonne & Intitulés acceptés & Oblig. & Valeur \\
\midrule\endhead
Code & Code ; Code EP & & 30 caractères au plus, unique dans l'activité \\
Intitulé & Intitulé ; Nom ; Name ; Title & \oblig & \\
Type & Type ; Code type ; Type code & \oblig & code d'un type d'EP de l'activité qui
  passe par le catalogue (voir la page des types) \\
ECTS & ECTS ; Crédits & \oblig & nombre strictement positif \\
Année min, max & Année min ; Année minimale ; Min year (et max) & & année d'étude, 0 pour
  aucune limite \\
Places & Places ; Capacité ; Capacity & & 0 pour un nombre illimité \\
État & État ; Etat ; State & & \col{publié} (ou \code{ouvert}), \code{clos} (ou
  \col{fermé}), \col{archivé}, \code{brouillon}, \col{proposé} ; en anglais
  \code{published}, \code{closed}, \code{archived}, \code{draft}, \code{proposed} \\
Inscriptions & Ouverture des inscriptions ; Fermeture des inscriptions & & dates, avec heure
  éventuelle \\
Déroulement & Début ; Fin & & dates \\
Responsables & Responsables ; Owners ; Teachers & & adresses ou numéros, séparés par des
  virgules ou des espaces ; ils s'ajoutent aux responsables existants \\
Description & Description & & texte \\
Évaluation & Évaluation ; Grading & & \code{type} (règle du type), \code{sans note},
  \code{avec note} (ou \code{note}) ; en anglais \code{none}, \code{graded} \\
Note maximale, seuil & Note maximale ; Seuil ; Seuil de validation & & nombres \\
Prérequis & Prérequis ; Prerequisites & & codes ou intitulés d'EP, séparés par des virgules
  ou des points-virgules ; un EP plus loin dans le fichier peut être cité \\
ECTS minimum & ECTS minimum ; Min ECTS & & crédits déjà acquis exigés à l'inscription \\
Prérequis bloquants & Prérequis bloquants & & oui/non \\
\bottomrule
\end{longtable}

\paragraph{Syllabus.} Les colonnes suivantes renseignent la fiche de l'EP (rubriques de la
fiche ECTS) : \code{Objectifs}, \code{Acquis d'apprentissage}, \col{Prérequis (texte)},
\code{Contenu}, \col{Méthodes}, \code{Conditions d'examen}, \code{Seconde session},
\code{Bibliographie}, \code{Langue}, \code{Travail personnel} (heures). Pour chaque EP dont au
moins une de ces cellules est remplie, une \textbf{nouvelle version du syllabus est publiée}
directement, sans passer par la revue : les cellules remplies remplacent le texte en vigueur,
les cellules vides le laissent tel quel, et les séances sont reprises. Une version soumise par
l'enseignant et non encore tranchée est reprise par l'import. Un import rejoué à l'identique
ne publie pas de nouvelle version.

\section{Séances des syllabus}

Même page, onglet «~séances~» (\code{import_catalog.php?kind=sessions}). Une ligne par
séance ; les séances d'un même EP se suivent dans l'ordre du fichier. Pour chaque EP cité, le
programme du syllabus en vigueur est \textbf{remplacé} par les séances du fichier, dans une
nouvelle version publiée qui reprend les autres rubriques.

\begin{longtable}{@{}L{2.6cm}L{4.6cm}L{1.4cm}L{6.6cm}@{}}
\toprule
Colonne & Intitulés acceptés & Oblig. & Valeur \\
\midrule\endhead
Code EP & Code ; Code EP & \oblig & code d'un EP du catalogue \\
Séance & Séance ; Intitulé ; Session ; Title & \oblig & \\
Nature & Nature ; Kind & & \code{CM} (cours magistral), \code{TD}, \code{TP}, \code{clinique},
  \code{terrain}, \code{projet}, \code{autre}, ou le libellé affiché par Moodle ; autre par
  défaut \\
Durée & Durée ; Duration ; Heures & & heures, décimales permises \\
Intervenant & Intervenant ; Lecturer ; Teacher & & adresse ou numéro d'un compte du site, sinon
  nom d'un intervenant extérieur gardé tel quel \\
\bottomrule
\end{longtable}

\section{Crédits}

Page \code{import_credits.php}, réservée à la DEVE (gestion de l'activité et
\code{mod/ep:validatedeve}) : elle porte des EP au dossier des étudiants, par exemple
l'historique d'un outil précédent ou des EP suivis dans un autre établissement. Les crédits
importés ont pour origine la DEVE ; un crédit tranché reçoit une décision dans son historique,
à la date indiquée.

\begin{longtable}{@{}L{2.8cm}L{5cm}L{1.4cm}L{6cm}@{}}
\toprule
Colonne & Intitulés acceptés & Oblig. & Valeur \\
\midrule\endhead
Étudiant & Étudiant ; Email ; Courriel ; Adresse de courriel ; Numéro d'identification ;
  Student & \oblig & adresse ou numéro d'un étudiant \emph{inscrit au cours} \\
Type & Type ; Code type & \oblig & code d'un type d'EP de l'activité, hors types de stage
  (leurs crédits viennent de la gestion des stages) \\
Code EP & Code EP ; EP du catalogue ; Activity code & & code d'un EP du catalogue, pour
  rattacher le crédit à cet EP \\
Intitulé & Intitulé ; Name ; Title & & obligatoire sans code EP (sinon, celui de l'EP) \\
Année & Année ; Année d'étude ; Study year & & 0 à 6 ; année courante de l'activité par
  défaut \\
ECTS demandés & ECTS demandés ; ECTS ; Claimed ECTS & \oblig & nombre \\
ECTS retenus & ECTS retenus ; Retained ECTS & & ECTS demandés par défaut \\
Note & Note ; Grade & & nombre \\
Statut & Statut ; Status & & \col{validé} (par défaut, ou \code{acquis}), \code{en attente},
  \col{refusé} (ou \col{non validé}), \col{annulé} (ou \col{retiré}) \\
Date de décision & Date de décision ; Date & & date ; jour de l'import par défaut \\
Validé par & Validé par ; Validateur ; Decided by & & adresse d'un compte du site ; l'auteur de
  l'import par défaut \\
Commentaire & Commentaire ; Comment & & \\
Description & Description & & \\
\bottomrule
\end{longtable}

Un crédit en attente entre dans le circuit de validation de son type, à la première étape.
Un crédit déjà porté au dossier (même étudiant, même type, même intitulé, même année, hors
crédits annulés) est écarté : un fichier importé deux fois ne double rien.

\section{Notes des participants d'un EP}

Rappel du lot~3 : depuis la liste des participants d'un EP du catalogue, le responsable ou la
DEVE importe les notes ou les validations (colonnes \code{Email} ou \col{Numéro
d'identification}, \code{Note} ou \col{Décision}, \code{Commentaire}). La liste exportée de
l'EP sert de modèle.

\section{Paramétrage (JSON)}

Page \code{settings_transfer.php}. Le paramétrage d'une activité s'exporte dans un fichier
JSON, s'importe dans une autre activité du même site ou d'un autre site, ou se copie
directement depuis une autre activité que l'on gère. Il ne contient \textbf{aucune donnée
d'étudiant}. On choisit les parties à importer :
\begin{itemize}
  \item \textbf{réglages généraux} : minimum de crédits du cursus, publication directe des
    mises à jour de syllabus, en-tête et pied des PDF, circuit et preuves du supplément,
    préfixe des numéros, champ de la date de naissance ;
  \item \textbf{types d'EP}, rapprochés par leur code (mis à jour ou créés) ;
  \item \textbf{minimums annuels} ;
  \item \textbf{supplément au diplôme} : catégories d'éléments (rapprochées par leur libellé)
    et textes des rubriques, en français et en anglais ;
  \item \textbf{modèles de courriels} ;
  \item \textbf{catalogue} : EP (rapprochés par code puis intitulé), responsables (par
    adresse), syllabus en vigueur et séances (intervenants par adresse), prérequis.
\end{itemize}
L'année courante n'est pas transférée : elle appartient à chaque promotion. Le fichier porte
un numéro de format ; un fichier d'un autre format est refusé plutôt que mal lu.

\section{Export complet et restauration}

Page \code{archive.php}. L'export complet (gestion de l'activité et consultation de tous les
dossiers) est une archive ZIP :
\begin{itemize}
  \item \code{classeur.xlsx} : classeur lisible dans un tableur, une feuille par sujet
    (bilan par étudiant, liste détaillée des crédits, catalogue, éléments du supplément,
    décisions d'année, suppléments émis). Les textes y sont toujours écrits comme du texte :
    une saisie qui commence par \code{=} ne devient pas une formule ;
  \item \code{parametrage.json} : le paramétrage, catalogue compris (section précédente) ;
  \item \code{donnees.json} : les données des étudiants -- crédits et l'historique de leurs
    décisions, éléments du supplément, décisions d'année, suppléments émis --, les personnes
    étant désignées par leur adresse et leur numéro d'identification ;
  \item \code{fichiers/} : justificatifs des crédits, preuves des éléments du supplément,
    suppléments émis (PDF certifiés).
\end{itemize}
La \textbf{restauration} (gestion de l'activité et validation DEVE) relit les deux fichiers
JSON, jamais le classeur : retoucher le classeur n'a pas d'effet. Elle ajoute à l'activité ce
que contient l'archive sans rien supprimer ; ce qui y figure déjà est écarté, si bien qu'une
archive restaurée deux fois ne double rien. Les étudiants sont retrouvés par leur numéro
d'identification puis leur adresse \emph{parmi les étudiants du cours} ; les données des
autres sont écartées et signalées. Un crédit d'origine «~stage~» devient un crédit saisi par
la DEVE (le stage d'origine appartient à l'autre activité). Un supplément émis garde son
empreinte de vérification ; si son numéro existe déjà sur le site, il reçoit un suffixe
\code{-R1}, \code{-R2}\ldots

Pour une sauvegarde de cours ou un transfert d'une année à l'autre, la sauvegarde Moodle du
cours reste l'outil de référence (elle comprend l'activité entière) ; l'export complet sert
à archiver une promotion hors de Moodle, à la remettre à un jury ou à reprendre des dossiers
dans une autre activité.
\end{document}
